\documentclass[12pt,letterpaper]{article}



%%
%% Required packages:
%%
\usepackage{etex}
\usepackage{amsmath}
\usepackage{amssymb}
\usepackage{amstext}
\usepackage{amsthm}
\usepackage{fancyhdr,fancybox}
\usepackage{mathrsfs} % need for \mathscr command
\usepackage{multicol}
\usepackage[normalem]{ulem}
%\usepackage{pictex,pictexplus}
% \usepackage{pictexwd}
\usepackage{rotating}
\usepackage{url}
\usepackage{xfrac}
\usepackage{booktabs}
\usepackage{color,soul}
\usepackage{comment}

\usepackage{hyperref}
\hypersetup{
    colorlinks=true,     % false: boxed links; true: colored links
    linkcolor=blue,      % color of internal links
    citecolor=blue,      % color of links to bibliography
    filecolor=blue,      % color of file links
    urlcolor=blue        % color of external links
}


\usepackage{tikz}
\usetikzlibrary{backgrounds,calc}

%%
%% Common formatting commands
%%
\setlength{\parindent}{0in}
\setlength{\textwidth}{7in}
\setlength{\evensidemargin}{-0.25in}
\setlength{\oddsidemargin}{-0.25in}
\setlength{\parskip}{.5\baselineskip}
\setlength{\topmargin}{-0.5in}
\setlength{\textheight}{9in}

%               Problem and Part
\newcounter{problemnumber}
\newcounter{partnumber}
\newcounter{subpartnumber}
\newcommand{\Problem}{%
  \refstepcounter{problemnumber}%
  \setcounter{partnumber}{0}%
  \setcounter{subpartnumber}{0}%
  \item[\makebox{\hfil\textbf{\theproblemnumber.}\hfil}]%
  }
\newcommand{\InProblem}[2][1.6]{%
  \refstepcounter{problemnumber}%
  \setcounter{partnumber}{0}%
  \setcounter{subpartnumber}{0}%
  \textbf{\theproblemnumber.}\ \ \parbox[t]{#1in}{#2}%
}
\newcommand{\InSmallProblem}[1]{\InProblem[1.05]{#1}}
\newcommand{\NoProblem}[1]{%
  \mbox{\hfil\textbf{#1}\hfil}%
  }
\newcommand{\UnProblem}{%
  \refstepcounter{problemnumber}%
  \setcounter{partnumber}{0}%
  \setcounter{subpartnumber}{0}%
  \mbox{\hfil\textbf{\theproblemnumber.}\hfil}%
  }
\newcommand{\InFlexProblem}[1]{%
  \refstepcounter{problemnumber}%
  \setcounter{partnumber}{0}%
  \setcounter{subpartnumber}{0}%
  \textbf{\theproblemnumber.}\ \ #1%
}
%%  Part:
\newcommand{\Part}{%
  \renewcommand{\thepartnumber}{(\alph{partnumber})}%
  \refstepcounter{partnumber}
  \setcounter{subpartnumber}{0}%
  \item[(\alph{partnumber})]%
}
\newcommand{\InPart}[2][1.75]{%
  \renewcommand{\thepartnumber}{(\alph{partnumber})}%
  \refstepcounter{partnumber}%
  \setcounter{subpartnumber}{0}%
  (\alph{partnumber})\ \ \parbox[t]{#1in}{#2}%
}
\newcommand{\UnPart}{%
  \refstepcounter{partnumber}%
  \renewcommand{\thepartnumber}{(\alph{partnumber})}%
  \setcounter{subpartnumber}{0}%
  (\alph{partnumber})%
}
\newcommand{\InLargePart}[1]{\InPart[3.05]{#1}}
\newcommand{\InFullPart}[1]{\InPart[6.2]{#1}}
\newcommand{\InSmallPart}[1]{\InPart[1.05]{#1}}
%% SubPart:
\newcommand{\SubPart}{%
  \refstepcounter{subpartnumber}%
  \item[(\roman{subpartnumber})]%
  }
\newcommand{\InSubPart}[2][2.25]{%
  \stepcounter{subpartnumber}%
  (\roman{subpartnumber})\ \ \parbox[t]{#1in}{#2}%
}
\newcommand{\InSmallSubPart}[1]{\InSubPart[1.5]{#1}}
\newcommand{\InTinySubPart}[1]{%
  \stepcounter{subpartnumber}%
  (\roman{subpartnumber})\ \ \makebox{#1}%
}
\newcommand{\InMySubPart}[2][2.25]{%
  \stepcounter{subpartnumber}%
  \parbox[t]{#1in}{\makebox[0.3in][l]{(\roman{subpartnumber})}\ \ {#2}}%
}
\newcommand{\TrueFalse}{\textbf{T}\ \ \textbf{F}\ \ }
\newcommand{\TFPart}[1]{% 
  \stepcounter{partnumber}%
  \setcounter{subpartnumber}{0}%
  \item[(\alph{partnumber})] \TrueFalse\parbox[t]{5.5in}{#1}}

\newcommand{\Points}[1]{(#1 points) \ }
\newcommand{\Point}[1]{(#1 point) \ }


%% For Answers:
\newcounter{answernumber}
\newcommand{\Answer}{%
  \refstepcounter{answernumber}%
  \setcounter{partnumber}{0}%
  \setcounter{subpartnumber}{0}%
  \item[\makebox{\hfil\textbf{\theanswernumber.}\hfil}]%
  }


% Setting up the headers for the first page
%\chead{} % will default to what is typed in the file.
\fancyhead[L]{\textsc{Math 22a}}
\fancyhead[R]{\textsc{Fall 2024}}
\fancyfoot[R]{
	\vspace*{\fill}\footnotesize\textcopyright{} \emph{Harvard Math 22a}	
}
\renewcommand{\headrulewidth}{0pt}

%\fancyhead[C]{}
%	\chead{}	
%	\lhead{}
%	\rhead{}
%	\cfoot{}


% Putting copyright on the footer of each page
%\fancypagestyle{secondpage}{
%	\fancyhead[C]{TESTing again} % change this to be blank
%		%\chead{}	
%	\fancyhead[L]{bears} % change this to be blank
%	\fancyhead[R]{dogs} % change this to be blank
%	\fancyfoot[R]{ %keeps the footer the same
%		\vspace*{\fill}\footnotesize\textcopyright{} \emph{Harvard Math 22a}	
%	}
%}
%\renewcommand{\headrulewidth}{0pt}
%\pagestyle{fancy}
\pagestyle{empty}


%\lhead{}
%\rhead{}
%\rfoot{\vspace*{\fill}\footnotesize\textcopyright{} \emph{Harvard Math 22a}}
%\renewcommand{\headrulewidth}{0pt}
%\pagestyle{fancy}




%%
%% Common shortcut commands:
%%
\newcommand{\ds}{\displaystyle}
\newcommand{\sgn}{\operatorname{sgn}}
\renewcommand{\Pr}{\operatorname{P}}
\newcommand{\CPr}[3][{}]{\Pr_{\text{#1}}\left(#2\,|\,#3\right)}

\newcommand{\C}{\mathbb{C}}
\newcommand{\R}{\mathbb{R}}
\newcommand{\Z}{\mathbb{Z}}
\newcommand{\N}{\mathbb{N}}
\newcommand{\Q}{\mathbb{Q}}
\newcommand{\D}{\mathbb{D}}

\newcommand{\rref}{\operatorname{rref}}
\newcommand{\im}{\operatorname{im}}
\newcommand{\rank}{\operatorname{rank}}
\newcommand{\diag}{\operatorname{diag}}
\newcommand{\Span}{\operatorname{span}}
%\renewcommand{\vec}[1]{\mathbf{#1}}
\newcommand{\charpoly}{\mathscr{P}}
\newcommand{\nicefrac}[2]{\sfrac{#1}{#2}} % using xfrac package
\newcommand{\topology}{\mathcal{T}}
\newcommand{\Inn}{\operatorname{Inn}}

\newcommand{\blankbox}[2]{\fbox{\rule{#1}{0in}\rule{0in}{#2}}}

\newenvironment{myindentpar}[1]%
 {\begin{list}{}%
         {\setlength{\leftmargin}{#1}
          \setlength{\rightmargin}{\leftmargin}}%
         \item[]%
 }
 {\end{list}}
\newtheorem{theorem}{Theorem}
\newtheorem*{NStheorem}{Theorem}
\newtheorem{cor}[theorem]{Corollary}
\newtheorem*{claim}{Claim}
\newtheorem{defn}{Definition}

\newenvironment{amatrix}[1]{%
  \left(\begin{array}{@{}*{#1}{c}|c@{}}
}{%
  \end{array}\right)
}

%--------grstep
% For denoting a Gauss' reduction step.
% Use as: \grstep{\rho_1+\rho_3} or \grstep[2\rho_5 \\ 3\rho_6]{\rho_1+\rho_3}
\newcommand{\grstep}[2][\relax]{%
   \ensuremath{\mathrel{
       {\mathop{\longrightarrow}\limits^{#2\mathstrut}_{
                                     \begin{subarray}{l} #1 \end{subarray}}}}}}
\newcommand{\swap}{\leftrightarrow}




\chead{\Large\textbf{Properties of Determinants, $\vec\S$3.1, 3.2}}% 

\begin{document}
\thispagestyle{fancy}

Recall from last class:

\begin{itemize}
\item For a matrix $A$ define  $A_{ij}$ to be the matrix obtained from $A$ by deleting row $i$ and column $j$.
\item For $n\geq 2$, the determinant of an $n\times n$ matrix $A=[a_{ij}]$ is
\begin{equation*}
\operatorname{det}(A)= a_{11} \operatorname{det}(A_{11}) - a_{12} \operatorname{det}(A_{12})+\dots + (-1)^{n+1} a_{1n} \operatorname{det}(A_{1n})=\sum_{j=1}^{n}(-1)^{1+j}a_{1j} \operatorname{det}(A_{1j}).
\end{equation*} We call this the cofactor expansion across the first row.
\item {\bf Mohan's Theorem:} Cofactor expansion across any row or column is the same, and hence each cofactor expansion gives the determinant.
\end{itemize}

\newpage




\begin{enumerate}

\Problem Compute the determinant of 
$C=\begin{bmatrix}
6 & 0 & 0 & 5\\
1 & 7 & 2 & -5\\
2 & 0 & 0 & 0\\
8 & 3 & 1 & 8\\
\end{bmatrix}$.

\vfill

\begin{defn} Let $A=[a_{ij}]$ be an $n\times n$ matrix, then $A$ is {\bf upper triangular} if $a_{ij}=0$ when $i>j$.
\end{defn}

{\bf Theorem 3.2:} Let $A$ be an $n\times n$ upper triangular matrix, then $$\operatorname{det}(A)= a_{11}a_{22}\dots a_{nn} = \prod_{j=1}^n a_{jj}.$$

\Problem Prove the theorem.

\vfill
\vfill


\newpage

{\bf Goal:} Prove $\operatorname{det}(A)\neq 0$ if and only if $A$ is invertible.

{\bf Theorem 3.3:} If $A$ is an $n\times n$ matrix and $E$ is an elementary matrix, then $$\operatorname{det}(EA)=\operatorname{det}(E) \operatorname{det}(A),$$ and 
$$\operatorname{det}(E)=\begin{cases}
1 & \text{if $E$ represents adding a multiple of a row to another row,}\\
-1 & \text{if $E$ represents interchanging two rows,}\\
r & \text{if $E$ represents scaling a row by a factor of $r$.}
\end{cases}$$

\Problem Let $$A=\begin{bmatrix} 1 & 2\\ 3 & 4\\ \end{bmatrix}, \quad E_1=\begin{bmatrix} 1 & 0\\ 2 & 1\\ \end{bmatrix}, \quad E_2=\begin{bmatrix} 0 & 1\\ 1 & 0\\ \end{bmatrix}, \quad E_3=\begin{bmatrix} 10 & 0\\ 0 & 1\\ \end{bmatrix}.$$ Check the theorem for these matrices.

\newpage

\Problem Prove the theorem.

\newpage

{\bf Theorem 3.6:} If $A$ and $B$ are $n$ by $n$ matrices, then $$\operatorname{det}(AB)=\operatorname{det}(A)\operatorname{det}(B).$$

\Problem Prove the theorem.

\vfill
\vfill

\Problem {\bf Corollary:} $A$ is invertible if and only if $\operatorname{det}(A)\neq0$. Prove the corollary.

\vfill

\newpage

\Problem {\bf True or False:} 

\Part If $A$ is an $n\times n$ invertible matrix, then $\operatorname{det}(A^{-1})=\frac{1}{\operatorname{det}(A)}$.

\vfill

\Part If $A$ is an $n\times n$ matrix, then $\operatorname{det}(A)=\operatorname{det}(A^{T})$.

\vfill

\Part If $A \sim B$, then $\operatorname{det}(A)=\operatorname{det}(B)$.

\vfill

\Problem We can think of $\operatorname{det}$ as a function from the set of $n\times n$ matrices to the real numbers. Is this function linear?

\vfill




\end{enumerate}


\begin{comment}
\newpage
\chead{\Large\textbf{Determinant Properties -- Answers / Solutions}}%
\lhead{}
\rhead{}
\thispagestyle{fancy}

\begin{enumerate}
  \Answer Again we use the definition to compute the determinant:
  \begin{align*}
  	\det C& =\det \begin{bmatrix}
  6 & 0 & 0 & 5\\
  1 & 7 & 2 & -5\\
  2 & 0 & 0 & 0\\
  8 & 3 & 1 & 8\\
  \end{bmatrix}= 6\cdot \det \begin{bmatrix}
  7 & 2 & -5\\
  0 & 0 & 0\\
  3 & 1 & 8\\
\end{bmatrix}-5\cdot \det \begin{bmatrix}
1 & 7 & 2\\
2 & 0 & 0\\
8 & 3 & 1\\
\end{bmatrix}=-5(-2)=10.
\end{align*}

Using the last Theorem, we can expand about the third row as well. We show that here:
  \begin{align*}
	\det C& =\det \begin{bmatrix}
		6 & 0 & 0 & 5\\
		1 & 7 & 2 & -5\\
		2 & 0 & 0 & 0\\
		8 & 3 & 1 & 8\\
	\end{bmatrix}= 2\cdot \det \begin{bmatrix}
		0 & 0 & 5\\
		7 & 2 & -5\\
		3 & 1 & 8\\
\end{bmatrix}
=2(5)(7-6)=10.
\end{align*}


  \Answer 
  
  \begin{proof}
  We prove the result by induction on the size of the matrix.
  
  {\bf Base Case:} Let $A=\begin{bmatrix} a & b\\ 0 & c\\ \end{bmatrix}$, then $\det(A)=ac$. Thus the result holds for $2\times 2$ matrices.
  
  {\bf Inductive Step:} We assume that the determinant of an $n\times n$ upper triangular matrix is given by the product of the diagonal elements.
  
  Let $A$ be an $(n+1)\times(n+1)$ upper triangular matrix. We compute the determinant by expanding along the first column:
  \begin{align*}
  \det(A)&= a_{11} \det(A_{11}) -a_{21}\det(A_{21})+\dots+(-1)^{n+1+1}a_{(n+1)1}\det(A_{(n+1)1})
  \end{align*}

Since $A$ is upper-triangular, $a_{21}=\dots=a_{(n+1)1}=0$. Thus
$$\det(A)= a_{11} \det(A_{11}).$$ Since $A_11$ is an $n\times n$ upper triangular matrix, we can apply the induction hypothesis to get $$\det(A)= a_{11} \det(A_{11})=a_{11} (a_{22}\dots a_{(n+1)(n+1)}).$$ Thus the result holds by mathematical induction.
  \end{proof}

  \Answer We first compute the products:
\begin{align*}
	E_1A &=\begin{bmatrix} 1 & 2\\ 5 & 8\\ \end{bmatrix},\\
	E_2A &= \begin{bmatrix} 3 & 4\\ 1 & 2\\ \end{bmatrix},\\
	E_3 A &= \begin{bmatrix} 10 & 20\\ 3 & 4\\ \end{bmatrix}.
\end{align*}
Thus notice
\begin{align*}
\det{E_1A}&= 8-10=-2 =\det{E_1}\det{A},\\
\det{E_2A}&=6-4 = 2 = \det{E_2}\det{A},\\
\det{E_3A}&=40-60=\det{E_3}\det{A}.
\end{align*}

Thus the Theorem seems to work in all these cases.
  
  \Answer   We prove the result by induction on the size of the matrix $A$.
  
  {\bf Base Cases:} Let $A=\begin{bmatrix} a & b\\ c & d\\ \end{bmatrix}$.
  
  {\bf Case 1:} Suppose $E$ corresponds to a $2$ by $2$ matrix scaling a row by a factor of $r$, then $E=\begin{bmatrix} r & 0\\ 0 & 1\\ \end{bmatrix}$ or $E=\begin{bmatrix} 1 & 0 \\ 0 & r\\ \end{bmatrix}$. Both matices are similar, so we only check the first. 
  
  Note $\det{EA}=\det \begin{bmatrix} r a & r b\\ c & d\\ \end{bmatrix}=rad-rbc=r(ad-bc)\det{E}\det{A}$, and hence the result holds in this case.
  
  {\bf Case 2:} Suppose $E$ corresponds to a $2$ by $2$ matrix interchaning rows, then $E=\begin{bmatrix} 0 & 1\\ 1 & 0\\ \end{bmatrix}$.
  
  Note $\det{EA}=\det \begin{bmatrix} c & d\\ a & b\\ \end{bmatrix}=bc-ad=-(ad-bc)=\det{E}\det{A}$, and hence the result holds in this case.
  
  {\bf Case 3:} Suppose $E$ corresponds to a $2$ by $2$ matrix adding $r$ times a row to another row, then $E=\begin{bmatrix} 1 & r\\ 0 & 1\\ \end{bmatrix}$ or $E=\begin{bmatrix} 1 & 0 \\ r & 1\\ \end{bmatrix}$. Both matices are similar, so we only check the first. 
  
  Note $\det{EA}=\det \begin{bmatrix} a+rc & b+rd\\ c & d\\ \end{bmatrix}=(a+rc)d-c(b+rd)=ad-bc=\det{E}\det{A}$, and hence the result holds in this case.
  
  
  {\bf Inductive Step:} Suppose the result holds for $n$ by $n$ matrices. Let $A$ be an $(n+1)\times (n+1)$ matrix and let $E$ be an elementary matrix. Since $n\geq 2$ and elementary row operations change at most 2 rows, we know there must be some row of $EA$ that is unchanged. We call that row $k$ and expand the determinant along that row. Let $C=EA$, then $C_{kj}$ is an $n$ by $n$ submatrix related to $A_{kj}$ by the same elementary row operation as $E$. By the induction hypothesis, we get $\det{C_{kj}}=\det{E} \det{A_{kj}}$ for all $1\leq j \leq n+1$. Now we compute the determinant of $EA$ expanding along row $k$, then
  \begin{align*}
  	\det{EA}=\det{C} &= (-1)^{k+1} a_{k1} \det{C_{k1}}+\dots+ (-1)^{n+1+k}a_{k(n+1)} \det{C_{k(n+1)}}\\
  	& = (-1)^{k+1} a_{k1} \det{E} \det{A_{k1}}+\dots+ (-1)^{n+1+k}a_{k(n+1)}\det{E} \det{A_{k(n+1)}}\\
  	&= \det{E} \left((-1)^{k+1} a_{k1} \det{A_{k1}}+\dots+ (-1)^{n+1+k}a_{k(n+1)} \det{A_{k(n+1)}} \right)\\
  	&=\det{E}\det{A}.
  \end{align*} 
Thus by mathematical induction the result holds for all $n$ by $n$ matrices.
    
  \Answer 
  
  \begin{proof}
  	
  We prove the result by cases.
  
  {\bf Case 1:} $A$ is invertible.
  
  Since $A$ is invertible, the invertible matrix theorem tells us that there exists elementary row operations $E_1, \dots, E_p$ such that $I_n=E_p \dots E_1 A$. Since elementary matrices are invertible and the inverse of an elementary matrix is another elementary matrix, we get $A=E_1^{-1}\dots E_p^{-1}$. Thus we get $AB=E_1^{-1}\dots E_p^{-1}B$. Now we apply the previous theorem multiple times to get the following
  \begin{align*}
  		\det{AB} &=\det{E_1^{-1}\dots E_p^{-1} B}\\
	&= \det{E_{1}^{-1}}\det{E_2^{-1}}\dots\det{E_{p}^{-1}}\det{B}\\
	&= \det{E_1^{-1}\dots E_p^{-1}}\det{B}\\
	&= \det{A}\det{B}.
\end{align*}
  
  {\bf Case 2:} $A$ is not invertible.
  
There must be a sequence of elementary matrices $E_1,\dots, E_p$ such that $A=E_1\dots E_p U$ where $U$ is the reduced row echelon form of $A$. Since $A$ is not invertible, the invertible matrix theorem gives that $U\neq I_n$. Moreover, $U$ must be upper triangular and have a zero along the diagonal. Thus $\det{A}=\det{E_1\dots E_p U}=\det{E_1}\dots \det{E_p}\det{U}=0$ since $\det{U}=0$.

Since $A$ is not invertible, we know $AB$ is not invertible. Thus we can apply the same reasoning to $AB$ to conclude $\det{AB}=0$. Hence $\det{AB}=0=\det{A}\det{B}$.
  \end{proof}
  
  \Answer 
  
  \begin{proof} We prove the two implications.
  	
  $\implies$ Suppose $A$ is invertible, then $AA^{-1}=I_n$. By the previous theorem, $1=\det{A}\det{A^{-1}}$, and hence $\det{A}\neq 0$.
  
  $\impliedby$ Suppose $\det{A}\neq 0$. We know there are elementary matrices $E_1,\dots E_p$ such that $A=E_1\dots E_p U$ where $U$ is the reduced row echelon form of $U$. We know $\det{E_k}\neq 0$ for all $1\leq k \leq p$. Thus $\det{U}\neq 0$. Since $U$ is an upper triangular matrix, we know the determinant of $U$ is the product of the diagonal elements. Thus $U$ cannot have a zero along the diagonal. Thus $U$ has all ones along the diagonal. Since $U$ is a square matrix, we get $U=I_n$. By the invertible matrix theorem, $A$ must be invertible.
  	
  \end{proof}
  
  \Answer 
  
  \begin{enumerate}
  \item {\bf True.} Since $A$ is invertible, we know $AA^{-1}=I_n$. Thus $1=\det{I_n}=\det{AA^{-1}}=\det{A}\det{A^{-1}}$, and hence $\det{A^{-1}}=\frac{1}{\det{A}}$.
  
  \item {\bf True.} Let $a_{ij}$ be the entry in row $i$ and column $j$ of $A$ and let $b_{ij}$ be the entry in row $i$ and column $j$ of $A^{T}$.
  
  Expand the determinant of $A$ along the first column and use $a_{ij}=b_{ji}$; namely, 
 \begin{align*}
 	\det{A}&=a_{11} \det A_{11} -a_{21}\det A_{21}+\dots+(-1)^{n+1} a_{(n+1)1}\det A_{(n+1)1}\\
 	&= b_{11} \det (A^{T})_{11}- b_{12} \det (A^{T})_{12}+\dots +(-1)^{n+1}b_{1n}\det{(A^{T})_{1n}}=\det{A^T}.
  \end{align*}
  
  \item {\bf False.} Let $A=\begin{bmatrix} 1 & 0\\ 0 & 1\\ \end{bmatrix}$ and $B=\begin{bmatrix} 0 & 1\\ 1 & 0\\ \end{bmatrix}$, then $A\sim B$ and $\det{A}=1$ and $\det{B}=-1$.
  \end{enumerate}
  
  \Answer Let $M_{n\times n}$ be the set of $n\times n$ matrices, then $\det : M_{n\times n} \to \mathbb{R}$ is a function. This function fails to be linear since $\det{I_2}=1$ but $\det{2I_2}=4$.
  
  It is worth noting that if we instead defined the function $f_j: \mathbb{R}^n \to \mathbb{R}$ by $\det{[\vec{a}_1 \dots, a_{j-1} \; \vec{x} \; \vec{a}_{j+1} \dots \vec{a}_n]}$, then $f_j$ is linear. Thus the determinant is linear in each column; we call this property {\bf multi-linear}.
  
\end{enumerate}

\end{comment}
\end{document}


%%% Local ables: 
%%% mode: latex
%%% TeX-master: t
%%% End: 
